\section{The cycle of nitrogen in the soil}
Nitrogen is the nutrient that limits the growth of most crops, and it is the one that the soil loses most easily. The nitrogen of the soil is in four forms. \emph{Organic nitrogen} (more than 95 per cent of the total) is in the residues, the humus and the microbes, and is not available to the plants. \emph{Ammonium} ($\mathrm{NH_4^+}$) is held by the clay and moves little. \emph{Nitrate} ($\mathrm{NO_3^-}$) is dissolved in the water, taken up by the crop and lost by leaching and denitrification. \emph{Gaseous forms} ($\mathrm N_2$, $\mathrm{N_2O}$, $\mathrm{NH_3}$) are lost to the air. The processes that move nitrogen from one form to another are:
\begin{exlist}
\item \kw{Mineralisation}: organic nitrogen $\to$ ammonium, by the microbes that decompose the organic matter.
\item \kw{Immobilisation}: mineral nitrogen $\to$ organic nitrogen, when the microbes take it to build their cells.
\item \kw{Nitrification}: ammonium $\to$ nitrate, in two steps, by aerobic bacteria; it needs oxygen, a pH above 5 and a warm soil.
\item \kw{Denitrification}: nitrate $\to$ $\mathrm{N_2O}$ $\to$ $\mathrm{N_2}$, by anaerobic bacteria in wet soil with organic carbon.
\item \kw{Volatilisation}: ammonium $\to$ ammonia gas, from the surface of an alkaline soil or of a spread manure.
\item \kw{Fixation}: $\mathrm N_2\to$ organic nitrogen, by the bacteria in the nodules of legumes.
\end{exlist}

\section{The carbon-to-nitrogen ratio}
The proportion of carbon to nitrogen in a material, the \kw{C/N ratio}, decides whether it gives or takes mineral nitrogen when it decomposes. The microbes have a C/N of about 8 and use about a third to a half of the carbon of the substrate for growth; they therefore need a substrate with a C/N below about 20 to 25 to find enough nitrogen. A residue with a lower ratio (young clover, 15) releases nitrogen; one with a higher ratio (straw, 80 to 100; sawdust, over 200) takes mineral nitrogen from the soil for weeks, and the crop that follows may suffer. This is the rule that the examples of this chapter apply.

\section{Factors that control the biological rates}
The rates of the processes of the soil depend on: \emph{temperature} (the rate doubles for about each ten degrees, the $Q_{10}$ rule, between 0 and 35 $^\circ$C); \emph{moisture} (a dry soil stops the microbes; a wet one lacks oxygen, so that nitrification stops and denitrification starts); \emph{aeration}; \emph{pH} (nitrification slows below pH 5.5); the \emph{supply of substrate} (the organic matter and its C/N); and the \emph{clay content}, which protects the organic matter from attack. In a model each factor is a number between 0 and 1 that multiplies the maximum rate.

\section{The movement of solutes in soil}
A dissolved substance moves by \kw{convection} (with the flowing water, at the velocity of the pore water), by \kw{diffusion} (from high to low concentration, in the water at rest) and by \kw{mechanical dispersion} (the spread caused by the different velocities in different pores); the last two are described by one dispersion coefficient. A substance that is held on the soil (\kw{sorption}) moves more slowly, and one that is taken up or decays is lost along the way. The equation of this chapter (Equation~\ref{eq:dm-solute}) is the balance of these terms. Nitrate does not sorb on most soils and moves with the water; ammonium and most pesticides sorb.

\section{Nitrogen and the water quality}
The nitrate that leaves the root zone reaches the groundwater and the surface waters. Nitrate above 50 mg\,L$^{-1}$ (11.3 mg N\,L$^{-1}$) is unfit for drinking water, and in the sea it causes the excess growth of algae (eutrophication). Nitrous oxide is a greenhouse gas about 270 times as strong as carbon dioxide over a century. The aim of management is to apply the nitrogen when the crop needs it, in the amount it needs, and to keep the soil covered in winter, when the drainage is largest. The nitrogen budget of Section~\ref{ch:map3-n} is the tool to find the surplus.

\section{Pesticides and other chemicals in soil}
A \kw{pesticide} is a chemical that kills a pest (a herbicide kills weeds, a fungicide fungi, an insecticide insects). Its fate in the soil is decided by sorption (a strongly sorbed chemical stays near the surface and is lost by erosion), degradation (the half-life, from days to years; faster in a warm, moist soil with an active biology), volatilisation, uptake by plants and leaching. The retardation factor and the half-life of Section~\ref{ch:map3-chem} are the two numbers that give a first estimate of the risk to the groundwater: a chemical that is weakly sorbed and slowly degraded is the most dangerous.
